% ═══════════════════════════════════════════════════════════════════════════════
% 中国大学生计算机设计大赛 — 人工智能实践赛作品报告
% 完全复刻 docs/作品报告模板.docx 格式
% 编译: xelatex report.tex
% ═══════════════════════════════════════════════════════════════════════════════

\documentclass[12pt,a4paper,linespread=1.5,fontset=none]{ctexart}

% ─── ctex 编号: 第一章、1.1、1.1.1（匹配 docx absNumId=0） ────────────────
\ctexset{
    section = {
        name = {第,章},
        number = \arabic{section},
    },
    subsection = {
        number = \arabic{section}.\arabic{subsection},
    },
    subsubsection = {
        number = \arabic{section}.\arabic{subsection}.\arabic{subsubsection},
    },
}

% ─── Packages ─────────────────────────────────────────────────────────────────
\usepackage[top=2.54cm, bottom=2.54cm, left=3.18cm, right=3.18cm]{geometry}
\usepackage{graphicx}
\usepackage{booktabs}
\usepackage{caption}
\usepackage{amsmath}
\usepackage{amssymb}
\usepackage{hyperref}
\usepackage{enumitem}
\usepackage{fancyhdr}
\usepackage{array}
\usepackage{multirow}
\usepackage{xcolor}
\usepackage{float}
\usepackage{fontspec}
\usepackage{titlesec}
\usepackage{indentfirst}

% ─── Fonts（匹配 docx: 宋体正文 / 黑体标题 / 华文中宋封面 / 华文楷体副标题）───
\setCJKmainfont{Songti SC}[BoldFont=Heiti SC]
\setCJKsansfont{Heiti SC}
\setCJKmonofont{STFangsong}
\setCJKfamilyfont{zhsong}{Songti SC}        % 宋体 → 正文
\setCJKfamilyfont{zhhei}{Heiti SC}           % 黑体 → 标题
\setCJKfamilyfont{zhkai}{Kaiti SC}           % 楷体 → 封面副标题
\setCJKfamilyfont{zhfs}{STFangsong}          % 仿宋
\setCJKfamilyfont{zhzongsong}{Songti SC}     % macOS 无华文中宋，宋体 Bold 替代

\newcommand{\songti}{\CJKfamily{zhsong}}
\newcommand{\heiti}{\CJKfamily{zhhei}}
\newcommand{\kaishu}{\CJKfamily{zhkai}}
\newcommand{\fangsong}{\CJKfamily{zhfs}}
\newcommand{\zongsong}{\CJKfamily{zhzongsong}}

% ─── Heading spacing（精确匹配 docx EMU 值）───────────────────────────────────
% H1: sb=46.8pt sa=7.8pt, 黑体 18pt Bold 居中
\titleformat{\section}{\centering\heiti\bfseries\fontsize{18pt}{27pt}\selectfont}{}{0em}{}
\titlespacing{\section}{0pt}{46.8pt}{7.8pt}

% H2: sb=23.4pt sa=4.7pt, 黑体 14pt Bold
\titleformat{\subsection}{\heiti\bfseries\fontsize{14pt}{21pt}\selectfont}{}{0em}{}
\titlespacing{\subsection}{0pt}{23.4pt}{4.7pt}

% H3: sb=15.6pt sa=3.1pt, 黑体 12pt Bold
\titleformat{\subsubsection}{\heiti\bfseries\fontsize{12pt}{18pt}\selectfont}{}{0em}{}
\titlespacing{\subsubsection}{0pt}{15.6pt}{3.1pt}

% ─── Body text ────────────────────────────────────────────────────────────────
% 正文段落: 宋体 12pt, 首行缩进 2 字符(≈24pt), 段前后 3.1pt
\setlength{\parindent}{2em}
\setlength{\parskip}{3.1pt plus1pt minus1pt}

% ─── Caption（图/表题注: sb=sa=4.7pt, 黑体 10pt）─────────────────────────────
\captionsetup[figure]{font={small,bf}, belowskip=4.7pt, aboveskip=4.7pt}
\captionsetup[table]{font={small,bf}, belowskip=4.7pt, aboveskip=4.7pt}

% ─── List spacing ─────────────────────────────────────────────────────────────
\setlist{nosep, leftmargin=2em, itemindent=0em}

% ─── Header ───────────────────────────────────────────────────────────────────
\pagestyle{fancy}
\fancyhf{}
\fancyhead[C]{\small\kaishu 基于改进 YOLOv8 的小目标交通标志实时检测}
\fancyfoot[C]{\thepage}
\renewcommand{\headrulewidth}{0.4pt}

\begin{document}

% ═══════════════════════════════════════════════════════════════════════════════
% 封面（完全复刻 docx）
% ═══════════════════════════════════════════════════════════════════════════════
\begin{titlepage}
\thispagestyle{empty}
\begin{center}

\vspace*{3cm}

% 华文中宋 28pt Bold 居中
{\zongsong\bfseries\fontsize{28pt}{42pt}\selectfont 2025年（第18届）\par}

\vspace{0.5cm}

% 华文中宋 28pt Bold 两端对齐（DISTRIBUTE）
{\zongsong\bfseries\fontsize{28pt}{42pt}\selectfont 中国大学生计算机设计大赛\par}

\vspace{1.5cm}

% 华文楷体 20pt 居中，段前后 31.2pt
\vspace{31.2pt}
{\kaishu\fontsize{20pt}{30pt}\selectfont 人工智能实践赛作品报告\par}
\vspace{31.2pt}

\vspace{2cm}

% 宋体 16pt 左对齐，段前后 7.8pt
\begin{flushleft}
{\songti\fontsize{16pt}{24pt}\selectfont
\vspace{7.8pt}
作品编号：\underline{\hspace{6cm}}\\[7.8pt]
\vspace{7.8pt}
作品名称：\underline{\hspace{6cm}}\\[7.8pt]
\vspace{7.8pt}
填写日期：\underline{\hspace{6cm}}\\[7.8pt]
}
\end{flushleft}

\end{center}
\end{titlepage}

% ═══════════════════════════════════════════════════════════════════════════════
% 第1章  作品概述
% ═══════════════════════════════════════════════════════════════════════════════
\section{作品概述}

随着智能交通系统的快速发展，交通标志检测成为自动驾驶与辅助驾驶领域的核心技术之一。然而在实际道路场景中，远处的小尺寸交通标志（在图像中占比 $< 0.1\%$，面积 $< 32\times 32$ 像素）往往难以被通用目标检测器有效识别，导致漏检率居高不下。

本作品针对小目标交通标志检测这一痛点，基于 YOLOv8s 检测框架，提出了一种融合多尺度检测、通道注意力机制与数据增强策略的改进方案。具体改进包括：(1) 增加 P2 高分辨率检测层，提升小目标特征表达能力；(2) 在 Neck 网络中引入高效通道注意力（ECA），增强关键通道响应；(3) 采用 Copy-Paste 数据增强，缓解小目标正样本稀疏问题；(4) 使用 Soft-NMS 后处理替代标准 NMS，减少密集场景下的误删除。

系统基于 TT100K 交通标志数据集（2021 版）进行训练与验证，选取 5 类典型小目标标志（限速 5、限速 30、限速 40、禁止驶入、禁止行人通行），最终在验证集上 mAP@0.5 达到 \textbf{0.844}，验证了改进方案的有效性。系统支持图片、视频及实时摄像头输入，可用于车载嵌入式设备或交通监控场景。

\textbf{关键词：}小目标检测；交通标志识别；YOLOv8；通道注意力；数据增强

% ═══════════════════════════════════════════════════════════════════════════════
% 第2章  问题分析
% ═══════════════════════════════════════════════════════════════════════════════
\section{问题分析}

\subsection{问题来源}

随着自动驾驶技术从 L2 向 L3/L4 级别演进，环境感知系统对交通标志检测的精度和实时性提出了越来越高的要求。在实际驾驶场景中，车辆需要在远距离（50--100 米）处提前识别交通标志，而此时标志在图像中往往仅占 $16\times 16$ 至 $32\times 32$ 像素区域，属于典型的小目标检测问题。

通用目标检测算法（如 YOLO、Faster R-CNN 等）在 COCO 等通用数据集上表现优异，但在小目标场景下性能下降明显。以 YOLOv8s 为例，其默认的三个检测头（P3/P4/P5）中，最高分辨率的 P3 层对应下采样倍数为 8，对于 $640\times 640$ 的输入图像，每个特征图像素对应 $8\times 8$ 的原始区域。一个 $16\times 16$ 的标志在该层上仅占据 $2\times 2$ 个特征点，极易被下采样过程中丢失的细节信息所淹没。

\subsection{现有解决方案}

针对小目标检测问题，学术界与工业界主要提出了以下几类解决方案：

\begin{enumerate}[label=(\arabic*)]
    \item \textbf{多尺度特征融合}：FPN \cite{lin2017fpn} 及 PANet、BiFPN 等通过自顶向下和自底向上的路径增强多尺度特征表达，但默认融合范围通常只覆盖 P3--P7，对更细粒度的 P2 特征利用不足。
    \item \textbf{注意力机制}：SE-Net、CBAM、ECA-Net \cite{wang2020eca} 等通过通道或空间注意力增强关键特征，其中 ECA 以极小的参数开销（1D 卷积替代全连接）实现了与 SE 相当的提升效果。
    \item \textbf{数据增强}：Mosaic、MixUp、Copy-Paste \cite{ghiasi2021simple} 等策略通过合成新样本增加训练数据多样性，Copy-Paste 特别适合小目标场景。
    \item \textbf{后处理优化}：Soft-NMS \cite{bodla2017soft} 以置信度衰减替代硬删除，在密集目标场景中能保留更多的真阳性预测。
\end{enumerate}

\subsection{本作品要解决的痛点问题}

\begin{enumerate}[label=(\arabic*)]
    \item \textbf{小目标特征丢失}：通用检测器的最高分辨率检测层（P3）对小目标而言特征粒度仍然不足；
    \item \textbf{特征通道冗余}：Neck 网络中所有通道等权处理，未考虑不同通道对小目标特征的贡献差异；
    \item \textbf{小目标正样本稀疏}：交通标志数据集中小目标实例有限，特别是尾部类别样本严重不足；
    \item \textbf{密集场景误删除}：城市道路中交通标志常成组出现，标准 NMS 容易误删邻近的低置信度框。
\end{enumerate}

% ═══════════════════════════════════════════════════════════════════════════════
% 第3章  技术方案
% ═══════════════════════════════════════════════════════════════════════════════
\section{技术方案}

\subsection{总体技术路线}

本作品以 YOLOv8s 为基线检测框架，采用四阶段递进式改进策略（图~\ref{fig:pipeline}），构建面向小目标交通标志的实时检测系统。四个改进阶段依次为：Stage 2 增加 P2 高分辨率检测层、Stage 3 引入 ECA 通道注意力、Stage 4 采用 Copy-Paste 数据增强、Stage 5 集成 Soft-NMS 后处理。

\begin{figure}[H]
\centering
\fbox{\begin{minipage}{0.85\textwidth}
\vspace{3cm}
\centering
\songti （此处插入系统总体技术路线框图）
\vspace{3cm}
\end{minipage}}
\caption{系统总体技术路线}
\label{fig:pipeline}
\end{figure}

\subsection{增加 P2 高分辨率检测层}

YOLOv8s 默认检测头包含 P3（下采样 8 倍，$80\times 80$）、P4（$40\times 40$）、P5（$20\times 20$）三层。本作品在 Neck 的 FPN 路径中增加一层上采样，将 P3 特征进一步放大至 P2 级别（下采样 4 倍，$160\times 160$），连接 Backbone 第 2 层（stride=4）经 C2f 融合后，构成 P2 检测头。P2 的每个空间位置对应 $4\times 4$ 像素区域，$16\times 16$ 的小目标在该层占据 $4\times 4$ 个特征点（P3 仅 $2\times 2$），特征表达能力翻倍。

改进后 P2 模型参数量由 11.1M 增至 14.5M（+30\%），FLOPs 由 28.7G 增至 47.5G（+65\%）。

\begin{figure}[H]
\centering
\fbox{\begin{minipage}{0.8\textwidth}
\vspace{2.5cm}
\centering
\songti （此处插入 P2 检测头结构示意图）
\vspace{2.5cm}
\end{minipage}}
\caption{P2 检测层结构示意（虚线框内为新增模块）}
\label{fig:p2}
\end{figure}

\subsection{高效通道注意力（ECA）}

在 Neck 网络的每个 C2f 模块后插入 ECA 模块。ECA 通过一维卷积捕获局部跨通道交互，参数量仅增加约 0.01M。其前向过程为：

\begin{equation}
y = \sigma\left(\text{Conv1D}_{k}\left(\text{GAP}(x)\right)\right) \cdot x
\end{equation}

其中卷积核大小 $k = |\frac{\log_2 C}{\gamma} + \frac{b}{\gamma}|_{\text{odd}}$（推荐 $\gamma=2, b=1$）。

\subsection{Copy-Paste 数据增强}

训练时以 0.5 概率执行：从同 batch 另一张图中随机选取小目标框（面积 $< 32^2$ px），复制像素并随机缩放 0.8--1.2 倍，粘贴到当前图的空旷区域（与现有框 IoU $< 0.1$），同步更新标注。该策略可有效增加尾部类别（如 p10 禁止行人通行）的正样本密度。

\subsection{Soft-NMS 后处理}

以高斯衰减替代标准 NMS 的硬删除，公式为：

\begin{equation}
s_i = s_i \cdot \exp\left(-\frac{\text{IoU}(M, b_i)^2}{\sigma}\right), \quad \forall b_i \notin D
\end{equation}

其中 $M$ 为当前最高分框，$\sigma=0.5$。在密集交通标志场景中可保留更多真阳性预测。

% ═══════════════════════════════════════════════════════════════════════════════
% 第4章  系统实现
% ═══════════════════════════════════════════════════════════════════════════════
\section{系统实现}

\subsection{数据集}

本作品采用 TT100K 2021 版数据集，原始数据包含约 100,000 张腾讯街景全景图，标注约 30,000 个交通标志实例，涵盖 221 类中国交通标志。筛选其中 5 类典型小目标标志进行实验，数据集类别分布如表~\ref{tab:dataset} 所示。

\begin{table}[H]
\centering
\caption{数据集类别分布}
\begin{tabular}{cccc}
\toprule
\textbf{类名} & \textbf{TT100K 标识} & \textbf{中文含义} & \textbf{标注数} \\
\midrule
speed\_limit\_5  & i2   & 限速 5 km/h      & 217  \\
speed\_limit\_30 & i4   & 限速 30 km/h     & 362  \\
speed\_limit\_40 & i5   & 限速 40 km/h     & 794  \\
no\_entry        & pne  & 禁止驶入         & 1120 \\
no\_pedestrians  & p10  & 禁止行人通行     & 193  \\
\bottomrule
\end{tabular}
\label{tab:dataset}
\end{table}

数据集按 7:2:1 划分为训练集（2171 张）、验证集（622 张）和测试集（314 张）。训练时下采样至 $640\times 640$，标注格式由原始 TT100K JSON 转换为 YOLO 归一化格式。

\subsection{实验环境}

\begin{table}[H]
\centering
\caption{实验环境配置}
\begin{tabular}{ll}
\toprule
\textbf{项目} & \textbf{配置} \\
\midrule
操作系统  & Ubuntu 22.04 / macOS 26.6 \\
GPU       & NVIDIA RTX 4060 (8GB) / Apple M5 \\
深度学习框架 & PyTorch 2.x + Ultralytics 8.x \\
Python    & 3.12 \\
\bottomrule
\end{tabular}
\label{tab:env}
\end{table}

\subsection{训练策略}

所有实验统一设置：输入分辨率 $640\times 640$，SGD 优化器（momentum=0.937，weight\_decay=0.0005），初始学习率 0.01，余弦退火调度，前 3 轮 warm-up，batch size 16（P2 模型降为 12），训练 100 轮，最后 10 轮关闭 Mosaic 增强。每阶段严格保持控制变量，仅变更目标改进点。

\subsection{消融实验设计}

为量化各改进点的独立贡献，设计消融实验矩阵如表~\ref{tab:ablation} 所示。

\begin{table}[H]
\centering
\caption{消融实验矩阵}
\begin{tabular}{ccccc}
\toprule
\textbf{实验} & \textbf{P2} & \textbf{ECA} & \textbf{Copy-Paste} & \textbf{Soft-NMS} \\
\midrule
A (基线)    & $\times$ & $\times$ & $\times$ & $\times$ \\
B           & $\checkmark$ & $\times$ & $\times$ & $\times$ \\
C           & $\checkmark$ & $\checkmark$ & $\times$ & $\times$ \\
D           & $\checkmark$ & $\checkmark$ & $\checkmark$ & $\times$ \\
E (完整)    & $\checkmark$ & $\checkmark$ & $\checkmark$ & $\checkmark$ \\
\bottomrule
\end{tabular}
\label{tab:ablation}
\end{table}

\subsection{系统功能模块}

\begin{itemize}
    \item \textbf{数据处理模块}（\texttt{prepare\_dataset.py}）：支持三种 JSON 格式自动检测与 YOLO 转换；
    \item \textbf{训练模块}（\texttt{train.py}）：YAML 配置驱动，CLI 参数覆写；
    \item \textbf{评估模块}（\texttt{eval.py}）：COCO 风格 mAP 评估，按小/中/大目标分组统计；
    \item \textbf{推理模块}（\texttt{demo.py}）：支持摄像头实时检测，显示 FPS。
\end{itemize}

% ═══════════════════════════════════════════════════════════════════════════════
% 第5章  测试分析
% ═══════════════════════════════════════════════════════════════════════════════
\section{测试分析}

\subsection{Stage 1：基线模型}

基线 YOLOv8s 在 TT100K 子集上训练 100 轮后的测试指标如表~\ref{tab:baseline} 所示。

\begin{table}[H]
\centering
\caption{基线模型测试指标}
\begin{tabular}{lc}
\toprule
\textbf{指标} & \textbf{数值} \\
\midrule
mAP@0.5      & 0.844 \\
mAP@0.5:0.95 & 0.589 \\
Precision    & 0.819 \\
Recall       & 0.807 \\
\bottomrule
\end{tabular}
\label{tab:baseline}
\end{table}

\subsection{Stage 2：P2 检测层}

增加 P2 检测层后，模型参数量由 11.1M 增至 14.5M（+30\%），FLOPs 由 28.7G 增至 47.5G（+65\%）。

（训练完成后填入具体指标与对比分析）

\subsection{消融实验汇总}

（各 Stage 训练完成后填入完整消融实验表格）

% ═══════════════════════════════════════════════════════════════════════════════
% 第6章  作品总结
% ═══════════════════════════════════════════════════════════════════════════════
\section{作品总结}

\subsection{作品特色与创新点}

\begin{enumerate}[label=(\arabic*)]
    \item \textbf{渐进式改进策略}：P2 检测层 $\to$ ECA 注意力 $\to$ Copy-Paste 增强 $\to$ Soft-NMS，四阶段递进，各阶段贡献可量化；
    \item \textbf{小目标专项优化}：P2 检测层将最高分辨率由 $80\times 80$ 提升至 $160\times 160$，小目标特征分辨率翻倍；
    \item \textbf{轻量化注意力}：ECA 仅增 0.01M 参数，兼顾精度与速度；
    \item \textbf{工程化设计}：YAML 配置驱动 + CLI 覆写，支持 Linux/Windows/Mac 一键训练，代码结构清晰可复现。
\end{enumerate}

\subsection{应用推广}

\begin{itemize}
    \item \textbf{车载辅助驾驶}：嵌入车载计算平台，为驾驶员提供实时交通标志提示；
    \item \textbf{交通监控}：部署于路边摄像头，自动检测并统计交通标志状态；
    \item \textbf{高精地图更新}：结合 SLAM 定位，批量采集道路标志信息用于地图更新。
\end{itemize}

\subsection{作品展望}

\begin{enumerate}[label=(\arabic*)]
    \item \textbf{模型轻量化}：探索知识蒸馏或模型剪枝，降低 P2 模型计算开销（当前 FLOPs 增加 65\%），达到边缘部署要求；
    \item \textbf{多天气鲁棒性}：引入雾天、夜间数据增强或域自适应方法，提升恶劣天气下的检测鲁棒性；
    \item \textbf{多帧融合}：利用视频时序信息进行目标跟踪，减少单帧漏检；
    \item \textbf{端到端部署}：导出 ONNX / TensorRT 模型，实现边缘设备上的实时推理。
\end{enumerate}

% ═══════════════════════════════════════════════════════════════════════════════
% 参考文献
% ═══════════════════════════════════════════════════════════════════════════════
\begin{thebibliography}{99}

\bibitem{lin2017fpn}
T.-Y. Lin, P. Doll\'{a}r, R. Girshick, K. He, B. Hariharan, and S. Belongie,
``Feature Pyramid Networks for Object Detection,''
in \textit{Proc. IEEE CVPR}, 2017, pp. 2117--2125.

\bibitem{wang2020eca}
Q. Wang, B. Wu, P. Zhu, P. Li, W. Zuo, and Q. Hu,
``ECA-Net: Efficient Channel Attention for Deep Convolutional Neural Networks,''
in \textit{Proc. IEEE CVPR}, 2020, pp. 11534--11542.

\bibitem{ghiasi2021simple}
G. Ghiasi, Y. Cui, A. Srinivas, R. Qian, T.-Y. Lin, E. D. Cubuk, Q. V. Le, and B. Zoph,
``Simple Copy-Paste is a Strong Data Augmentation Method for Instance Segmentation,''
in \textit{Proc. IEEE CVPR}, 2021, pp. 2918--2928.

\bibitem{bodla2017soft}
N. Bodla, B. Singh, R. Chellappa, and L. S. Davis,
``Soft-NMS --- Improving Object Detection with One Line of Code,''
in \textit{Proc. IEEE ICCV}, 2017, pp. 5561--5569.

\bibitem{yolov8}
G. Jocher, A. Chaurasia, and J. Qiu,
``Ultralytics YOLOv8,''
2023. [Online]. Available: \url{https://github.com/ultralytics/ultralytics}

\bibitem{tt100k}
Z. Zhu, D. Liang, S. Zhang, X. Huang, B. Li, and S. Hu,
``Traffic-Sign Detection and Classification in the Wild,''
in \textit{Proc. IEEE CVPR}, 2016, pp. 2110--2118.

\end{thebibliography}

\end{document}
